\newcommand{\avg}[1]{\left< #1 \right>} % for average
\let\underdot=\d % rename builtin command \d{} to \underdot{}
\renewcommand{\d}[2]{\frac{d #1}{d #2}} % for derivatives
\newcommand\given[1][]{\:#1\vert\:}
\newcommand{\expected}[1]{\left\langle #1 \right\rangle}
\newcommand{\kseq}{\bm{k}}
\newcommand{\dir}[1]{#1^{\scriptscriptstyle\updownarrow}}
\newcommand{\untri}[1]{\tilde{#1}} % the independent "half" of the symmetric matrices of undir case
\newcommand{\abs}[1]{\left| #1 \right|}
\newcommand{\ull}[1]{\underline{#1}}
\renewcommand{\epsilon}{\varepsilon}
\newcommand{\eps}{\varepsilon}

%%%%%%%%%%%%%%%%%%%%%%%%%%%%%%
%% ESTILOS DE TEOREMAS Y ENTORNOS
%%%%%%%%%%%%%%%%%%%%%%%%%%%%%%
\theoremstyle{plain}
\newtheorem{theorem}{Theorem}[chapter] % Crea el contador 'theorem'
\newtheorem{lemma}[theorem]{Lemma}     % CORREGIDO: usa 'theorem' en lugar de 'thm'
\newtheorem{cor}[theorem]{Corollary}   % CORREGIDO: comparte contador con theorem
\crefname{cor}{Corollary}{Corollaries}
\newtheorem{conj}[theorem]{Conjecture} % CORREGIDO: usa 'theorem' en lugar de 'thm'

\theoremstyle{definition}
\newtheorem{definition}[theorem]{Definition} % Ahora comparte la misma secuencia numérica
\newtheorem{example}[theorem]{Example}       % AÑADIDO: Esto salva el Capítulo 5

%%%%%%%%%%%%%%%%%%%%%%%%%%%%%%
%% COLORES Y OTROS ESTILOS
%%%%%%%%%%%%%%%%%%%%%%%%%%%%%%
\definecolor{realnet}{rgb}{0.9,0.9,0.85}
\definecolor{unbiasnet}{HTML}{00BFC4}
\definecolor{biasnet}{HTML}{F8766D}

\newtheoremstyle{question}{6pt}{6pt}{\bfseries}{}{\bfseries}{:}{ }{}
\theoremstyle{question}
\newtheorem{qq}{\sc RQ}[]
\crefname{qq}{Research question}{research questions}

%%%%%%%%%%%%%%%%%%%%%%%%%%%%%%
%% COMANDOS MATEMÁTICOS
%%%%%%%%%%%%%%%%%%%%%%%%%%%%%%
\newcommand{\review}[1]{{\color{blue} #1}}
\newcommand{\N}{\mathbb{N}}
\newcommand{\R}{\mathbb{R}}
\newcommand{\floor}[1]{\left \lfloor{#1}\right \rfloor }
\DeclareRobustCommand{\mean}[1]{\left\langle #1 \right\rangle}
\newcommand{\RA}{\Rightarrow}
%\renewcommand{\G}{\mathcal{G}}

\newcommand{\AIC}{\text{AIC}}
\newcommand{\BIC}{\text{BIC}}

\newcommand{\MHD}{\text{MHD}}

\newcommand{\model}[1]{#1}

\newcommand{\vecb}[1]{\boldsymbol{#1}}
\DeclareMathOperator*{\argmin}{arg\,min}
\DeclareMathOperator*{\argmax}{arg\,max}
\DeclareMathOperator{\E}{\mathbb{E}}
\DeclareMathOperator{\sign}{sign}
